% ============================================================================
% templates/exam/solucionario/solucionario.tex — plantilla: solucionario de examen
% ----------------------------------------------------------------------------
% Para: transformar un examen existente de la colección en un solucionario
% formal, con portada, enunciado íntegro y desarrollo completo.
% El modo [soluciones] muestra los desarrollos y suprime los espacios de
% respuesta, de modo que el MISMO archivo sirve como examen y solucionario.
% Uso: scripts/build.sh solucionario.tex   (LuaLaTeX)
% ============================================================================
\documentclass[soluciones]{academic-exam}

\universidad{Universidad Nacional de San Cristóbal de Huamanga}
\facultad{Facultad de Ciencias Económicas, Administrativas y Contables}
\escuela{Escuela Profesional de Economía}
\curso{Economía Internacional II}
\codigocurso{EC-454}
\docente{Mg. Nombre Apellido}
\periodo{Semestre 2021-I}
\tipoevaluacion{Examen Parcial 02}
\fechaevaluacion{12 de junio de 2021}
\duracion{100 minutos}
\modalidad{Presencial}
\autorevaluacion{Edison Achalma}

\begin{document}
\portada[Expediente: \texttt{economia-internacional-ii/parciales/2021-1-parcial-02}]

\begin{observacion}[Sobre este solucionario]
Los enunciados se transcriben literalmente del examen original; los
desarrollos son elaboración propia y pueden diferir del solucionario
oficial del docente. Cuando existan varios caminos de solución se indica
el más directo y se comenta la alternativa.
\end{observacion}


\begin{pregunta}[5][Tipo de cambio real]
El tipo de cambio nominal es $E = 3{,}80$ soles por dólar, el nivel de
precios de EE.\,UU.\ es $P^{*} = 110$ y el nivel de precios interno es
$P = 104{,}5$.
\begin{apartados}
  \item Calcule el tipo de cambio real y explique su significado. \pts{2}
  \item Si en un año la inflación interna es 6\,\%, la externa 2\,\% y la
        depreciación nominal 3\,\%, determine qué ocurre con la
        competitividad. \pts{3}
\end{apartados}
\begin{solucion}
\textbf{a)} El tipo de cambio real es
\[
  q = \frac{E\,P^{*}}{P} = \frac{3{,}80 \times 110}{104{,}5} = 4{,}00:
\]
una canasta estadounidense cuesta el equivalente a 4 canastas... en
rigor, $q$ indica cuántas canastas internas se requieren para adquirir
una canasta externa; valores mayores implican un país más
\emph{competitivo} (producción interna relativamente barata).

\textbf{b)} En tasas: $\hat q = \hat E + \pi^{*} - \pi
= 3\,\% + 2\,\% - 6\,\% = -1\,\%$. El tipo de cambio real se
\emph{aprecia} 1\,\%: la depreciación nominal no compensa el diferencial
de inflación y la competitividad se deteriora.
\end{solucion}
\end{pregunta}

\begin{pregunta}[5][Balanza de pagos]
Clasifique cada operación en la cuenta corriente o en la cuenta
financiera de la balanza de pagos peruana, indicando el signo:
\begin{apartados}
  \item Exportación de arándanos por USD~40 millones.
  \item Utilidades reinvertidas por una minera de capital extranjero.
  \item Compra de bonos soberanos peruanos por un fondo chileno.
  \item Remesas recibidas de trabajadores peruanos en el exterior.
\end{apartados}
\begin{solucion}
\textbf{a)} Cuenta corriente, crédito ($+$): exportación de bienes.
\textbf{b)} Doble registro: renta de la inversión extranjera en cuenta
corriente ($-$, débito) y reinversión como flujo de inversión directa en
la cuenta financiera ($+$).
\textbf{c)} Cuenta financiera, pasivo/entrada de cartera ($+$).
\textbf{d)} Cuenta corriente, ingreso secundario, crédito ($+$).
\end{solucion}
\end{pregunta}

\begin{pregunta}[10][Modelo Mundell-Fleming]
Considere una economía pequeña y abierta con tipo de cambio flexible y
perfecta movilidad de capitales ($r = r^{*}$).
\begin{apartados}
  \item Represente el equilibrio inicial en el plano $(Y, r)$ con las
        curvas IS*, LM* y la recta $r = r^{*}$. \pts{2}
  \item Analice gráficamente y explique el efecto de una expansión
        fiscal sobre el producto, el tipo de cambio y la balanza
        comercial. \pts{4}
  \item Contraste con el resultado bajo tipo de cambio fijo y concluya
        sobre la efectividad relativa de la política fiscal en cada
        régimen. \pts{4}
\end{apartados}
\begin{solucion}
\textbf{a)} Equilibrio donde IS y LM se cortan sobre la recta
$r = r^{*}$; con $r > r^{*}$ entran capitales (presión a la apreciación)
y viceversa.

\textbf{b)} La expansión fiscal desplaza IS a la derecha y presiona $r$
al alza; la entrada de capitales aprecia el sol; la apreciación reduce
exportaciones netas y devuelve la IS a su posición inicial. Resultado:
$\Delta Y = 0$, apreciación del tipo de cambio y deterioro de la balanza
comercial exactamente igual al impulso fiscal (\emph{crowding out}
externo completo).

\textbf{c)} Con tipo de cambio fijo, la presión a la apreciación obliga
al banco central a comprar divisas, expandiendo la oferta monetaria: la
LM acompaña a la IS y el producto aumenta con toda la potencia del
multiplicador. Conclusión clásica: la política fiscal es \emph{ineficaz}
con cambio flexible y \emph{plenamente eficaz} con cambio fijo; la
monetaria, al revés (trilema de la economía abierta).
\end{solucion}
\end{pregunta}

\begin{observacion}[Comentario final]
Errores frecuentes en este examen: confundir la definición de $q$ (usar
$P/EP^{*}$ e invertir la lectura de competitividad) y omitir el doble
registro de las utilidades reinvertidas en la pregunta 2. Conviene
repasar la convención de signos del manual de balanza de pagos del FMI
(sexta edición).
\end{observacion}

\findelexamen
\end{document}
